\section{Экспериментальная проверка и оценка результатов}

\subsection{Методика испытаний}

Проверка системы проводилась на локальном стенде, поднимаемом из \texttt{deploy/docker/docker-compose.yml}. Перед началом серии тестов фиксировались версии образов (Master, Agent, PostgreSQL, HAProxy, Loki, Grafana) и параметры переменных окружения, влияющих на диапазон TCP-портов и пути динамической конфигурации HAProxy.

\textbf{Сценарии.} Набор проверок согласован с каталогом функциональных сценариев в приложении~Б: регистрация и heartbeat агента, публикация и жизненный цикл compose-проекта, добавление и удаление TCP-маршрутов (в том числе с автоматическим listen), перезагрузка HAProxy после изменения файла, отказ при конфликте порта, просмотр сгенерированного конфига в UI. Для каждого сценария фиксировались: входные данные, ожидаемый результат по критериям ниже, фактический статус HTTP и наблюдения в логах Loki.

\textbf{Контроль среды.} Отдельно проверялась согласованность сетевой модели Docker: доступность \texttt{host.docker.internal} из контейнера HAProxy к порту, опубликованному compose на хосте агента. На Linux-хостах без Docker Desktop эквивалентная схема может потребовать иной адресации backend --- это зафиксировано как ограничение стенда, а не как дефект кода при условии корректной настройки.

\textbf{Воспроизводимость.} Каждый ключевой сценарий выполнялся минимум дважды: «холодный'' старт после \texttt{docker compose up} и «тёплый'' повтор с уже заполненной БД и существующими маршрутами. Это выявляет ошибки инициализации и ошибки, связанные с накопленным состоянием.

\subsection{Критерии успешности}

Сценарий считается успешным, если выполняются условия:
\begin{itemize}[leftmargin=1.5cm]
  \item API возвращает консистентный статус и понятные диагностические сообщения;
  \item в UI отображается фактическое состояние маршрутов;
  \item сгенерированный \texttt{chronos-tcp.cfg} проходит валидацию HAProxy;
  \item TCP-подключения по listen-порту доходят до целевого backend-порта;
  \item перезапуск компонентов не приводит к потере критического состояния.
\end{itemize}

\subsection{Результаты функциональных испытаний}

В результате испытаний подтверждено:
\begin{itemize}[leftmargin=1.5cm]
  \item корректное добавление маршрутов с автоматическим выбором listen-порта в допустимом диапазоне;
  \item корректное удаление маршрутов и обновление конфигурации прокси;
  \item стабильная работа сценария проксирования к \texttt{host.docker.internal:PORT};
  \item успешное отображение маршрутов и их состояния в UI.
\end{itemize}

\textbf{Соответствие каталогу сценариев.} Сценарии F-01--F-04 (регистрация, heartbeat, публикация, запуск) и F-07--F-14 (TCP-маршруты, reload HAProxy, доступ к backend, просмотр cfg) на стенде автора выполнены успешно. Сценарии F-15--F-18 (песочница, очистка агентов, аудит, повторный подъём) проверялись выборочно: подтверждена работоспособность путей, критичных для демонстрации дипломного прототипа; полный регрессионный прогон рекомендуется вынести в отдельный контур CI.

Выявленные и устраненные в процессе реализации дефекты:
\begin{itemize}[leftmargin=1.5cm]
  \item ошибка HAProxy при наличии UTF-8 BOM в начале динамического конфига;
  \item проблемы запуска shell-скрипта перезагрузки при CRLF-окончаниях строк;
  \item неоднозначность пользовательской интерпретации backend-портов (порт контейнера vs порт хоста).
\end{itemize}

\subsection{Ограничения эксперимента}

Испытания носят функциональный характер: не измерялись предельная пропускная способность TCP-прокси, задержки при одновременной публикации десятков проектов и деградация БД при длительном накоплении аудита. Нагрузочное тестирование и формализованный отчёт SLA не входили в объём работы и относятся к направлениям развития.

\subsection{Обсуждение результатов}

Полученные результаты подтверждают рабочую гипотезу: для существенной доли практических задач возможно построить инженерно пригодную оркестрацию на основе Compose без перехода на тяжелые кластерные платформы. Ключевым фактором является наличие централизующего слоя (Master + UI) и стандартизованного агентного контура.

Система показала устойчивость в условиях частых изменений маршрутов и конфигураций. Наиболее критичные риски концентрируются в эксплуатационной плоскости:
\begin{itemize}[leftmargin=1.5cm]
  \item ошибки в сетевой адресации backend;
  \item неконсистентность окружения (различия Docker Desktop и Linux host-mode);
  \item неверная трактовка опубликованных портов.
\end{itemize}

\subsection{Направления дальнейшего развития}

В рамках следующих итераций целесообразно реализовать:
\begin{enumerate}[leftmargin=1.5cm]
  \item multi-master режим с отказоустойчивым лидерством;
  \item расширенную модель авторизации и ролей;
  \item автоматический discovery опубликованных портов compose-проектов агента;
  \item формализованные нагрузочные бенчмарки и тесты деградации;
  \item расширенную интеграцию с внешними системами наблюдаемости.
\end{enumerate}
